\documentclass[11pt]{article}

\newcommand{\myname}{Nirav Bhattad}
\newcommand{\myrollnumber}{23B3307}
\newcommand{\topicname}{Part-2}

\usepackage{nirav}
% \geometry{a4paper, margin=1in}

% \theoremstyle{definition}
% \newtheorem{definition}{Definition}

% Custom environment for Speaker Notes
\newmdenv[
  linecolor=blue,
  linewidth=2pt,
  backgroundcolor=blue!5,
  skipabove=10pt,
  skipbelow=10pt,
  innerleftmargin=10pt,
  innerrightmargin=10pt,
  innertopmargin=10pt,
  innerbottommargin=10pt
]{speakernote}

\setlength{\headheight}{14pt}

\title{Part 2: The Somewhat Homomorphic Scheme (SHE) \\ \large Speaker Notes \& Mathematical Primer}
\author{Nirav Bhattad}
\date{March 27, 2026}

\begin{document}

\maketitle

\section{The Symmetric Foundation}

\begin{speakernote}
\textbf{How to start speaking:} 
``We are now going to build a Fully Homomorphic Encryption scheme from scratch. To do this, we won't start with the complex public-key version. Instead, we are going to look at a symmetric-key version[cite: 17]. It is beautifully simple, and if you understand this symmetric version, the rest of the paper will fall perfectly into place.''
\end{speakernote}

Before defining the algorithms, we establish the parameters based on a security parameter $\lambda$:
\begin{itemize}
    \item $\eta$: The bit-length of the secret key.
    \item $\rho$: The bit-length of the noise.
\end{itemize}

The symmetric scheme is defined by three algorithms:

\begin{itemize}
    \item $\mathsf{KeyGen}(\lambda)$: Choose a random odd integer $p$ from the interval $[2^{\eta-1}, 2^\eta)$. The integer $p$ is the secret key.
    
    \item $\mathsf{Encrypt}(p, m)$: To encrypt a bit $m \in \{0, 1\}$, we choose a large random integer $q$ and a small random integer $r$ (where $|r| < 2^\rho$). We compute the ciphertext $c$ as:
    \begin{equation}
        c = p \cdot q + 2r + m
    \end{equation}
    
    \item $\mathsf{Decrypt}(p, c)$: To recover the message, we compute:
    \begin{equation}
        m' = (c \bmod p) \bmod 2
    \end{equation}
\end{itemize}

\begin{speakernote}
\textbf{Walking through the Decryption:}
``Let's pause and look at that decryption equation, because it is the engine of the entire paper. Our ciphertext is $c = pq + 2r + m$. What happens when we take this modulo $p$? 

Because $pq$ is an exact multiple of $p$, it vanishes. We are left with the remainder: $2r + m$. But wait—the modulo operation only returns $2r + m$ if this noise term is strictly less than $p/2$ in absolute value! If the noise grows too large, it ``wraps around'' the modulus $p$, and decryption fails. 

Assuming the noise is small enough, $(c \bmod p)$ gives us exactly $2r + m$. Now, what happens when we take $(2r + m) \bmod 2$? The $2r$ is an even number, so it vanishes, leaving only our original plaintext bit $m$. Simple, elegant, and completely reliant on the noise staying small.''
\end{speakernote}

\section{Transitioning to a Public-Key Scheme}

A symmetric scheme is useful, but a true FHE scheme must allow anyone to encrypt data We transform our symmetric scheme into a public-key scheme using a technique that relies on publishing a list of ``encryptions of zero.''

We introduce two new parameters:
\begin{itemize}
    \item $\gamma$: The bit-length of the integers in the public key.
    \item $\tau$: The number of integers in the public key.
\end{itemize}

\begin{itemize}
    \item $\mathsf{KeyGen}(\lambda)$: The secret key remains the odd integer $p$. To create the public key, we generate $\tau + 1$ samples of the form $x_i = p \cdot q_i + 2r_i$. We ensure $x_0$ is the largest element, $x_0$ is odd, and its noise term $r_0$ is even. The public key is $pk = \langle x_0, x_1, \dots, x_\tau \rangle$.
    
    \item $\mathsf{Encrypt}(pk, m)$: To encrypt a bit $m \in \{0, 1\}$, we choose a random subset $S \subseteq \{1, 2, \dots, \tau\}$ and a small random noise integer $r$. The ciphertext is:
    \begin{equation}
        c = \left[ m + 2r + 2\sum_{i \in S} x_i \right]_{x_0}
    \end{equation}
\end{itemize}

\begin{speakernote}
\textbf{Intuition for the Public Key:}
``Why does this work? Every $x_i$ in the public key is essentially an encryption of zero. By summing a random subset of these $x_i$ values, we are adding $0 + 0 + 0 \dots$ together, generating a brand new, randomized encryption of zero. We then add our message bit $m$ to it. The modulo $x_0$ step at the end is just to keep the physical size of the ciphertext from growing infinitely large, without destroying the underlying modulo $p$ structure.''
\end{speakernote}

\section{Homomorphism and The Noise Bottleneck}

We now have a working public-key encryption scheme. But is it homomorphic? To answer this, we must define the $\mathsf{Evaluate}$ algorithm. In this scheme, $\mathsf{Evaluate}$ is beautifully trivial: it simply applies integer addition and multiplication directly to the ciphertexts.

Let's observe what happens when we operate on two ciphertexts, $c_1 = p q_1 + 2r_1 + m_1$ and $c_2 = p q_2 + 2r_2 + m_2$.

\subsection{Homomorphic Addition}
When we add the two ciphertexts, we group the terms:
\begin{equation}
    c_1 + c_2 = p(q_1 + q_2) + 2(r_1 + r_2) + (m_1 + m_2)
\end{equation}
Notice the structure: we have a multiple of $p$, plus $2 \times (\text{new noise})$, plus the sum of the messages. The new noise is exactly $(r_1 + r_2)$. By the triangle inequality, addition only increases the noise linearly.

\subsection{Homomorphic Multiplication}
Now, let us multiply the ciphertexts:
\begin{align}
    c_1 \cdot c_2 &= (p q_1 + 2r_1 + m_1)(p q_2 + 2r_2 + m_2) \\
    &= p(p q_1 q_2 + q_1(2r_2 + m_2) + q_2(2r_1 + m_1)) \nonumber \\
    &\quad + 2(2r_1 r_2 + r_1 m_2 + r_2 m_1) + m_1 m_2
\end{align}
Again, the structure is preserved! We have a massive multiple of $p$, plus 2 times a mixed noise term, plus the product of the messages $m_1 m_2$. 

\begin{speakernote}
\textbf{The Climax of Part 2: The Noise Bottleneck}
``Look closely at the noise term in the multiplication equation: $2r_1 r_2$. When we multiply two ciphertexts, the noise doesn't just add; it multiplies. It squares the magnitude of the noise. 

Remember our strict decryption rule? $(c \bmod p)$ only works if the accumulated noise is strictly less than $p/2$. Because multiplication squares the noise, we can only perform a very limited number of sequential multiplications before the noise term explodes past $p/2$, destroying the plaintext. 

This is why the scheme is currently only \textbf{Somewhat Homomorphic}. It can evaluate any circuit, but only up to a very specific multiplicative depth. In Part 4, we will solve this using Gentry's Bootstrapping Theorem.''
\end{speakernote}

\end{document}